\section*{Розділ V. Порядок роботи ЦВКс}
\addcontentsline{toc}{section}{Розділ V. Порядок роботи ЦВКс}

\subsection*{5.1. Засідання ЦВКс}
\addcontentsline{toc}{subsection}{5.1. Засідання ЦВКс}
    5.1.1. ЦВКс працює у режимі засідань, які скликаються за потребою для вирішення питань, віднесених до її компетенції.

    5.1.2. Засідання ЦВКс скликаються:

        \begin{enumerate}[label=\alph*)]
            \item За ініціативою Голови ЦВКс;
            \item За вимогою не менше половини членів ЦВКс;
            \item За дорученням КСУ для організації конкретних виборів;
            \item У зв'язку з надходженням скарг або заяв, що потребують колегіального розгляду.
        \end{enumerate}

    5.1.3. Голова ЦВКс зобов'язаний скликати засідання не пізніше ніж через 3 робочі дні з дня надходження вимоги або доручення.

    5.1.4. Під час активних виборчих кампаній ЦВКс може засідати щоденно за потребою для оперативного вирішення виборчих питань.

\subsection*{5.2. Порядок скликання засідань}
\addcontentsline{toc}{subsection}{5.2. Порядок скликання засідань}
    5.2.1. Повідомлення про засідання ЦВКс надсилається всім членам комісії не пізніше ніж за 24 години до його проведення.

    5.2.2. У невідкладних випадках засідання може бути скликано без дотримання 24-годинного строку попередження за умови згоди всіх членів ЦВКс.

    5.2.3. Повідомлення про засідання повинно містити:

        \begin{enumerate}[label=\alph*)]
            \item Дату, час та спосіб проведення засідання;
            \item Проєкт порядку денного;
            \item Матеріали до питань порядку денного (за наявності);
            \item Технічну інформацію для участі у засіданні (для дистанційних засідань).
        \end{enumerate}

\subsection*{5.3. Форми проведення засідань}
\addcontentsline{toc}{subsection}{5.3. Форми проведення засідань}
    5.3.1. Засідання ЦВКс може проводитися у таких формах:

        \begin{enumerate}[label=\alph*)]
            \item Очно -- з фізичною присутністю членів ЦВКс у визначеному місці;
            \item Дистанційно -- з використанням засобів електронного зв'язку (відеоконференції, системи електронного голосування тощо);
            \item У змішаному форматі -- з можливістю участі як очно, так і дистанційно.
        \end{enumerate}

    5.3.2. Форма проведення засідання визначається Головою ЦВКс з урахуванням характеру питань, технічних можливостей та обставин.

    5.3.3. При проведенні засідання ЦВКс дистанційно або у змішаному форматі забезпечується:

        \begin{enumerate}[label=\alph*)]
            \item Належна ідентифікація учасників засідання;
            \item Можливість вільного обговорення питань порядку денного;
            \item Достовірна фіксація результатів голосування;
            \item Конфіденційність обговорення, якщо це необхідно;
            \item Ведення протоколу засідання.
        \end{enumerate}

\subsection*{5.4. Порядок проведення засідань}
\addcontentsline{toc}{subsection}{5.4. Порядок проведення засідань}
    5.4.1. Засідання ЦВКс відкриває та веде Голова ЦВКс, а у разі його відсутності -- заступник Голови.

    5.4.2. На початку засідання встановлюється кворум та затверджується порядок денний.

    5.4.3. Кожне питання порядку денного обговорюється колегіально з наданням можливості висловитися всім членам ЦВКс.

    5.4.4. За підсумками обговорення приймається відповідне рішення згідно з процедурами, визначеними пунктом 4.5 цього Положення.

    5.4.5. Засідання ЦВКс може бути закритим для сторонніх осіб у разі розгляду конфіденційних питань або персональних справ.

\subsection*{5.5. Документообіг та електронний документообіг}
\addcontentsline{toc}{subsection}{5.5. Документообіг та електронний документообіг}
    5.5.1. Документообіг ЦВКс ведеться в електронній формі з використанням засобів електронного цифрового підпису відповідно до законодавства України.

    5.5.2. Офіційні документи ЦВКс підписуються електронним цифровим підписом Голови ЦВКс та скріплюються електронним цифровим підписом Секретаря ЦВКс.

    5.5.3. Електронні документи ЦВКс мають таку ж юридичну силу, як і документи на паперових носіях, за умови дотримання вимог законодавства про електронний цифровий підпис.

    5.5.4. ЦВКс веде електронний архів усіх документів, пов'язаних з виборчим процесом, включаючи:

        \begin{enumerate}[label=\alph*)]
            \item Протоколи засідань ЦВКс;
            \item Рішення ЦВКс;
            \item Документи кандидатів та виборців;
            \item Виборчі протоколи та підсумкові документи;
            \item Скарги та заяви з виборчих питань;
            \item Переписку з органами студентського самоврядування.
        \end{enumerate}

    5.5.5. Доступ до електронного архіву забезпечується всім членам ЦВКс, а також може надаватися іншим особам відповідно до принципів прозорості та у межах, не заборонених законодавством.

\subsection*{5.6. Протоколювання засідань}
\addcontentsline{toc}{subsection}{5.6. Протоколювання засідань}
    5.6.1. Всі засідання ЦВКс протоколюються Секретарем ЦВКс.

    5.6.2. Протокол засідання ЦВКс повинен містити:

        \begin{enumerate}[label=\alph*)]
            \item Дату, час та місце (спосіб) проведення засідання;
            \item Список присутніх членів ЦВКс;
            \item Порядок денний засідання;
            \item Короткий виклад обговорення кожного питання;
            \item Результати голосування з зазначенням позицій кожного члена ЦВКс;
            \item Прийняті рішення;
            \item Особливі думки членів ЦВКс (за їх бажанням).
        \end{enumerate}

    5.6.3. Протокол засідання підписується електронним цифровим підписом Голови та Секретаря ЦВКс протягом 2 робочих днів після проведення засідання.

    5.6.4. Копії протоколів засідань надсилаються всім членам ЦВКс та можуть оприлюднюватися на офіційних інформаційних ресурсах органів студентського самоврядування (за винятком конфіденційної інформації).

\subsection*{5.7. Планування роботи ЦВКс}
\addcontentsline{toc}{subsection}{5.7. Планування роботи ЦВКс}
    5.7.1. ЦВКс на початку свого терміну повноважень розробляє план роботи на рік, який включає:

        \begin{enumerate}[label=\alph*)]
            \item Календар планованих виборів до органів студентського самоврядування;
            \item Строки підготовчих заходів до кожних виборів;
            \item Плани підвищення кваліфікації членів ЦВКс;
            \item Заходи з інформування студентської спільноти про виборчі процедури.
        \end{enumerate}

    5.7.2. План роботи ЦВКс подається до відома КСУ та може коригуватися протягом року за рішенням ЦВКс.

\subsection*{5.8. Взаємодія з технічними службами}
\addcontentsline{toc}{subsection}{5.8. Взаємодія з технічними службами}
    5.8.1. Для забезпечення технічної сторони виборчого процесу ЦВКс може залучати технічних працівників Університету або зовнішніх фахівців.

    5.8.2. Координацію роботи з технічними службами здійснює Секретар ЦВКс під загальним керівництвом Голови ЦВКс.

    5.8.3. Усі особи, залучені до технічного забезпечення виборчого процесу, зобов'язані дотримуватися принципів неупередженості та конфіденційності.